\section{Related Work}

\textbf{LLM Benchmarks.}
We briefly review the common LLM benchmarks, following the classification presented in \autoref{fig:eval_classification}. The most prevalent benchmarks are static, ground-truth-based ones, typically in the form of multiple-choice questions or question-answering tasks with predefined answers and test cases.
These benchmarks encompass a range of topics including language understanding, mathematics, coding, and logical reasoning. Prominent examples in this category are MMLU~\cite{hendrycks2020measuring}, HellaSwag~\cite{zellers2019hellaswag}, GSM-8K~\cite{cobbe2021gsm}, BigBench~\cite{srivastava2023beyond}, AGIEval~\cite{zhong2023agieval}, and HumanEval~\cite{chen2021evaluating}. Benchmarks focusing on safety, such as ToxicChat~\cite{toxicchat}, and comprehensive suites like HELM~\cite{liang2022holistic}, also exist.
In addition to closed-ended questions, benchmarks can include open-ended questions that are evaluated by human judgment, which can be rated by experts or crowd workers such as Amazon Mechanical Turk~\cite{karpinska-etal-2021-perils,koala_blogpost_2023,wang-etal-2023-self-instruct}.
The recent trend includes utilizing GPT-4 for approximating human judgment~\cite{chiang-lee-2023-large}, with notable instances being MT-Bench~\cite{zheng2023judging} and AlpacaEval~\cite{alpaca_eval}.
In addition to static benchmarks, live benchmarks that include fresh questions are also available. These questions can be obtained from annual exams or weekly online contests such as Codeforces~\cite{li2022competition, huang2023competition}. They can also be sourced from human interaction. Some studies have explored using live human interaction for reinforcement learning from human preference~\cite{bai2022training,ouyang2022training,touvron2023llama2}. However, these studies are typically limited to specific organizations.
In this paper, we introduce \system, the first open, large-scale, and crowdsourced benchmark platform that utilizes live human interaction.

\textbf{Risks of Static Benchmarks.}
Static benchmarks have certain issues, including contamination, saturation, overfitting, and a lack of human alignment~\cite{yang2023rethinking,oren2023proving}. DynaBench~\cite{kiela2021dynabench} identifies these challenges and recommends the use of a live benchmark that incorporates a human-in-the-loop approach for classical NLP benchmarks. Our system adopts a similar spirit. However, our focus is on chatting with LLMs, and we implement this on a significantly larger user scale.

\textbf{Ranking System.}
Ranking systems have been a well-studied topic in statistics. Related topics include probability models~\cite{mm_bradley_terry,tie_in_bradley_terry}, rank elicitation~\cite{online_rank,pmlr-v32-busa-fekete14,preference_rank_elicitation}, and online experiment design~\cite{sequential_design, karimi2021online}.
The Elo rating system has also been used for LLMs~\cite{bai2022training,boubdir2023elo}.
Contributing to this literature, we introduce techniques for accelerating ranking convergence and detecting abnormalities, specifically applied to large-scale, real-world settings of LLMs.

\textbf{Human Preference Dataset.}
Owing to the significance of human preferences, several datasets and analyses exist that incorporate human preferences. These include OpenAssistant~\cite{kopf2023openassistant}, HH-RLHF~\cite{bai2022training}, LMSYS-Chat-1M~\cite{zheng2023lmsyschat1m}, and synthetic approximations of human preferences like UltraFeedback~\cite{cui2023ultrafeedback} and Nectar~\cite{starling2023}.
Our prior data release, LMSYS-Chat-1M~\cite{zheng2023lmsyschat1m}, is similarly collected via crowdsourcing. However, LMSYS-Chat-1M comprises solely conversations and lacks human preference data, rendering it unsuitable for direct use in ranking studies. This paper focuses on the analysis of preference data for ranking purposes.


\iffalse

\section{Related Work}

\textbf{LLM Benchmarks}
We briefly introduce the common LLM benchmarks, according to the classification in \autoref{fig:eval_classification}.

- static ground-truth-based benchmark
MMLU~\cite{hendrycks2020measuring}, HellaSwag~\cite{zellers2019hellaswag}, GSM-8K~\cite{cobbe2021gsm},
BigBench~\cite{srivastava2023beyond}, AGIEval~\cite{zhong2023agieval}, Drop~\cite{dua2019drop}, and HumanEval~\cite{chen2021evaluating}. There are also safety-related benchmarks like ToxicChat~\cite{toxicchat} and comprehensive suites like HELM~\cite{liang2022holistic}.

- static benchmark for human preference
Super-NaturalInstructions~\cite{wang2022super}, MT-Bench~\cite{zheng2023judging}, AlpacaEval~\cite{alpaca_eval}, ~\cite{chiang-lee-2023-large},
Amazon Mechanical Turk~\cite{karpinska-etal-2021-perils,koala_blogpost_2023}.

- live benchmark  with ground truth
Codeforce~\cite{huang2023competition,li2022competition}.

- live benchmark with human preference
private ones at Anthropic~\cite{bai2022training}, Meta~\cite{touvron2023llama2}.
but we need an open one.

- other interactive benchmarks
AgentBench~\cite{liu2023agentbench}, WebArena~\cite{zhou2023webarena}, SWE-Bench~\cite{jimenez2024swebench}.

\textbf{Risks and solutions of static benchmark}
DynaBench~\cite{kiela2021dynabench}
Contamination~\cite{yang2023rethinking,oren2023proving,deng2023investigating}
ImageNet~\cite{recht2019imagenet}

\textbf{Data collection and analysis}
OpenAssistant~\cite{kopf2023openassistant}, Anthropic~\cite{bai2022training}, UltraFeedback~\cite{cui2023ultrafeedback},
LMSYS-Chat-1M~\cite{zheng2023lmsyschat1m},
Self-Instruct~\cite{wang-etal-2023-self-instruct}

Rank Elicitation~\cite{online_rank,pmlr-v32-busa-fekete14,preference_rank_elicitation},
ranking models~\cite{mm_bradley_terry,tie_in_bradley_terry},
online experiment design~\cite{sequential_design, karimi2021online}

\fi